\section{Discussion and Limits}

This paper now makes two claims: a conceptual one and a small executable one. Conceptually, we argue that serialization is a first-class design axis that should be reasoned about explicitly whenever irregular 3D structures must be aligned with efficient model families. Empirically, we show on a deterministic synthetic probe that locality-aware serializations expose local mixed-channel neighborhoods much more gracefully than raster or random baselines under tight sequence budgets. That second claim is intentionally narrower than a biomolecular benchmark claim.

We do not present new real-world biomolecular benchmark results, and we do not argue that serialized representations are universally better than native geometric ones. In some tasks, explicit geometry-aware operators or equivariant methods will remain the better choice. The stronger position is that serialization quality should be treated as something measurable and task-dependent rather than as a preprocessing afterthought.

This perspective also does not imply that every historical precursor deserves revival. The reason the source paper is worth upgrading is that later literature independently converged on related ideas for efficiency and scale. That downstream convergence gives retrospective significance to the original Hilbert-mapping formulation.

There are three practical risks in this area. The first is \emph{order-induced artifact learning}: models may learn regularities that reflect the traversal more than the underlying structure. The second is \emph{false-neighborhood accumulation}: if ordering repeatedly separates functionally linked regions or binds irrelevant ones together, scale alone will not rescue the representation. The third is \emph{synthetic-overreach}: a clean executable probe can still be too abstract if it is mistaken for a real scientific benchmark. These risks reinforce the importance of explicit evaluation criteria for locality, reversibility, channel semantics, and benchmark realism.

The publication implication is therefore clear. A strong follow-on empirical paper should not simply report accuracy under a new backbone. It should explain what representation contract the serialization creates, what scientific information survives that contract, and why the chosen ordering is appropriate for the task. The next genuinely stronger version of this line of work would port the same diagnostics to real protein structures, real atomistic channels, and real learned sequence backbones.
